% !TEX root = ../main.tex
\vspace{-2mm}
\section{Introduction}
\vspace{-1mm}
  \begin{figure}[ht!]
    \centering
    \includegraphics[width=1.05\linewidth, trim = 0.9cm 0.5cm 4.4cm 0.4cm, clip]{images/teaser.pdf}
    \caption{An overview of Diffusion Language Models (DLMs) (top) and the proposed \textbf{Inverse-distilled Diffusion Language Models (IDLM)} (bottom). DLMs train a denoising network $\widehat{f}$ using samples from the target data distribution $p^*(x_0)$. In contrast, our approach leverages inverse distillation to learn
    $p_{\theta}(x_0)$ such that the resulting solution recovers the pretrained model $f^*.$
    Consequently, while standard DLMs exhibit strong generative performance, they typically require
    % incur high inference latency due to 
    the large number of sampling steps. IDLM instead trains a few-step generator that maintains generative quality while substantially accelerating inference.}
    \vspace{-7mm}
  \label{fig:teaser}
  \end{figure}

  % Generative modeling for discrete data, such as natural text and code,
  % has gained widespread adoption as the output of large language models~\citep{touvron2023llama, groeneveld2024olmo, lozhkov2024starcoder}
  % becomes increasingly indistinguishable from real-world human data in terms of perplexity and quality~\citep{sadasivan2023can, kushnareva2024ai}.
  % Recently, Diffusion Language Models (DLMs; \citealp{sohl2015deep, austin2021structured, campbell2022continuous, lou2024discrete})
  % have emerged as promising framework for modeling the distribution of text data. General framework for broad class of DLMs was established by SEDD~\citep{lou2024discrete} model. However, the general case of DLMs is generally intractable in practice and leads to overcomplicated theoretical part. To overcome this issue, further works~\citep{sahoo2024simple, sahoo2025the, schiff2024discreteguidance, shi2024simplified} concentrated on specific types of processes and as result simplified theoretical formulation of DLMs. It leads to the fact that DLM achieved comparable performance in the field of text generation in comparison with dominant paradigm autoregressive (AR) language models~\citep{nie2025large}. In contrast to AR models DLMs offer the benefits of flexible, parallel generation. However, the iterative nature of the reverse diffusion process necessitates hundreds or thousands of sampling steps,
  % creating a significant latency bottleneck that hinders practical use.
  
Generative modeling for discrete data, such as natural language, has seen widespread adoption, driven by the rapid advancement of large language models~\citep{touvron2023llama, groeneveld2024olmo, lozhkov2024starcoder}. 
% The outputs of these models have become increasingly indistinguishable from human-generated content in terms of both perplexity and perceptual quality~\citep{sadasivan2023can, kushnareva2024ai}. 
Recently, Diffusion Language Models 
(DLMs; \citealp{sohl2015deep, austin2021structured, campbell2022continuous, lou2024discrete}) have emerged as a promising framework for modeling the distribution of the text data. A general formulation for a broad class of DLMs was introduced by~\citep{lou2024discrete}. However, the general formulation is often intractable in practice and introduces significant theoretical complexity. To address this, subsequent work~\citep{sahoo2024simple, sahoo2025the, schiff2024discreteguidance, shi2024simplified} has focused on specific types of diffusion processes, enabling simplified theoretical analyses and more tractable implementations. As a result, DLMs have achieved competitive performance in text generation, comparable to that of autoregressive language models~\citep{nie2025large}. However, the inherently iterative nature of reverse diffusion requires hundreds or even thousands of sampling steps at inference time, resulting in considerable latency that limits their practicality.

On the other hand, in the continuous domain, the issue of slow inference in diffusion models~\citep{ho2020denoising, song2020score, karras2022elucidating} has been the focus of extensive research~\citep{song2020denoising, luhman2021knowledge, zhang2022fast, lu2022dpm}. Among the various approaches, \emph{distillation} has emerged as one of the most effective solutions~\citep{song2023consistency, salimans2022progressive, luhman2021knowledge, berthelot2023tract, xie2024distillation}. However, directly applying such techniques in the discrete domain is challenging due to the non-differentiable nature of discrete data. Recent efforts have begun to address this limitation for DLMs. Notably, SDTT~\citep{deschenaux2024beyond} and Duo-DCD~\citep{sahoo2025the} propose adaptations of consistency-based distillation methods~\citep{song2023consistency, song2023improved, kim2023consistency} to the discrete setting, demonstrating initial progress toward enabling fast inference in DLMs.
  % By leveraging continuous-time Markov
  % chains to corrupt data and subsequently learning to reverse this noise process, architectures like SEDD~\citep{lou2024discrete}
  % and MDLM~\citep{sahoo2024simple} have demonstrated perplexity scores competitive with state-of-the-art autoregressive language models while offering the benefits of flexible, parallel generation.
  

  % On the other side in the domain of continuous diffusion models~\citep{ho2020denoising, song2020score, karras2022elucidating} the longstanding problem of slow inference trying to solve it for a long time~\citep{song2020denoising, luhman2021knowledge, zhang2022fast, lu2022dpm}. most notable example among all of these approaches is distillation~\citep{song2023consistency, salimans2022progressive, luhman2021knowledge, berthelot2023tract, xie2024distillation}. However, practically it is very hard to directly apply these technique due to the non-differentiable nature of the data. Recent works have begun
  % to address this for DLMs: SDTT~\citep{deschenaux2024beyond} and Duo-DCD~\citep{sahoo2025the}
  % employs adaption distillation technique, known as consistency models. 
  
  % In our work we employ another direction of distillation. Specifically, we extend the theoretical framework of \textit{Inverse Distillation}~\citep{gushchin2025inverse,kornilov2025universal} for continuous models to the discrete setting. This yields \textit{Inverse-distilled Diffusion Language Models (IDLM)}, a principled framework for accelerating a broad class of discrete diffusions.
  % Our main idea, is instead of training the student directly to mimic the teacher's outputs, we treat the pre-trained teacher as the fixed optimal solution to a denoising objective. We then optimize the student distribution to be the one for which this teacher is optimal.  
  In this work, we investigate an alternative distillation paradigm by extending the theoretical framework of \emph{Inverse Distillation}~\citep{gushchin2025inverse, kornilov2025universal}, originally developed for continuous diffusion models, to the discrete domain. However, this direct extension introduces several nontrivial challenges arising from the intrinsic properties of discrete data. First, \underline{theoretical concerns} emerge regarding the validity of the training objective, particularly due to the potential non-uniqueness of its minimizers. Second, \underline{practical challenges} arise from the non-smoothness of the discrete space, which complicates the optimization. 
  
  \textbf{Contributions.} To address the above-mentioned issues, we present \emph{Inverse-distilled Diffusion Language Models (IDLM)}, a principled framework for accelerating a broad class of discrete diffusion models. In summary, our core contributions are as follows:
  % The core principle of this approach is to treat the pretrained teacher model as a fixed optimal solution to a denoising objective. The student distribution is then optimized to be the one for which this teacher is optimal.
  % \textbf{Our main contribution} is the first method, \textit{IDLM}, that in a principled manner extends the \textit{inverse distillation} framework to \textit{discrete diffusion models}. In this regard, we:
  \vspace{-0.9em}

  \begin{enumerate}
    \setlength{\itemsep}{1pt}
    \setlength{\parskip}{1pt}
    \item We theoretically establish the uniqueness of the global optimum for the IDLM objective, validating the correctness of the proposed optimization procedure (\wasyparagraph\ref{sec:theoretical extension idlm}).

    \item We propose gradient-stable relaxations that enable effective training through (i) a \emph{simplex-parameterized} generator and (ii) a principled mechanism for \emph{backpropagating gradients through the discrete forward process} (\wasyparagraph\ref{sec:simplex relaxation}).
    \item We empirically demonstrate that IDLM achieves competitive \emph{generation quality}, matching the performance of 1024-step teacher models while requiring up to $64\times$ fewer sampling steps (\wasyparagraph\ref{sec:exps}).
  \end{enumerate}
